
\newcounter{paranum} 
\newcommand{\np}{\stepcounter{paranum}\noindent\textbf{\arabic{paranum}.}\quad}

\section{Concepts}

\np faithful, injective, projective, flat, Noetherian, semisimple, Artinian, Soc, essnetial epi
idem, Jacobson radical, projective cover, DVR.

Composition series, Wedderburn’s Theorem, idempotents, Schur's lemma, the cartan matrix,Krull-Schmidt, Socle Series,Radical Series,Top, Nakayama lemma,Idempotent Lifting
Jordan-Hölder, Loewy 

\section{Module}

\begin{theorem}[Universal property of direct sums] \label{thm:universal-property-direct-sum} 
    Let \(A\) be a unital associative ring, and let \(\{M_i\}_{i\in I}\) be a family of left \(A\)-modules. The direct sum \[ \bigoplus_{i\in I} M_i \] together with the canonical inclusions \[ \iota_i:M_i\longrightarrow \bigoplus_{i\in I}M_i \] satisfies the following universal property: for every left \(A\)-module \(N\) and every family of homomorphisms \(f_i:M_i\to N\), there exists a unique homomorphism \[ f:\bigoplus_{i\in I}M_i\longrightarrow N \] such that \(f\circ \iota_i=f_i\) for all \(i\in I\). Equivalently, \[ \operatorname{Hom}_A\left(\bigoplus_{i\in I}M_i,N\right) \cong \prod_{i\in I}\operatorname{Hom}_A(M_i,N). \] 
\end{theorem} 
    
\begin{theorem}[Universal property of direct products] \label{thm:universal-property-direct-product} 
    Let \(A\) be a unital associative ring, and let \(\{M_i\}_{i\in I}\) be a family of left \(A\)-modules. The direct product \[ \prod_{i\in I} M_i \] together with the canonical projections \[ \pi_i:\prod_{i\in I}M_i\longrightarrow M_i \] satisfies the following universal property: for every left \(A\)-module \(N\) and every family of homomorphisms \(g_i:N\to M_i\), there exists a unique homomorphism \[ g:N\longrightarrow \prod_{i\in I}M_i \] such that \(\pi_i\circ g=g_i\) for all \(i\in I\). Equivalently, \[ \operatorname{Hom}_A\left(N,\prod_{i\in I}M_i\right) \cong \prod_{i\in I}\operatorname{Hom}_A(N,M_i). \] 
\end{theorem}


\begin{theorem}[Tensor--Hom adjunction]
\label{thm:tensor-hom-adjunction}
Let \(A\) and \(B\) be unital associative rings.

\begin{enumerate}
    \item[(i)] Let \({}_B X_A\) be a \((B,A)\)-bimodule. Then the functor
    \[
    X\otimes_A -: A\text{-}\mathrm{Mod}\longrightarrow B\text{-}\mathrm{Mod}
    \]
    is left adjoint to
    \[
    \operatorname{Hom}_B(X,-):B\text{-}\mathrm{Mod}
    \longrightarrow A\text{-}\mathrm{Mod}.
    \]
    Equivalently, for every left \(A\)-module \(M\) and every left
    \(B\)-module \(N\), there is a natural isomorphism
    \[
    \operatorname{Hom}_B(X\otimes_A M,N)
    \cong
    \operatorname{Hom}_A\bigl(M,\operatorname{Hom}_B(X,N)\bigr).
    \]

    \item[(ii)] Let \({}_A X_B\) be an \((A,B)\)-bimodule. Then the functor
    \[
    -\otimes_A X:\mathrm{Mod}\text{-}A\longrightarrow \mathrm{Mod}\text{-}B
    \]
    is left adjoint to
    \[
    \operatorname{Hom}_B(X,-):\mathrm{Mod}\text{-}B
    \longrightarrow \mathrm{Mod}\text{-}A.
    \]
    Equivalently, for every right \(A\)-module \(M\) and every right
    \(B\)-module \(N\), there is a natural isomorphism
    \[
    \operatorname{Hom}_B(M\otimes_A X,N)
    \cong
    \operatorname{Hom}_A\bigl(M,\operatorname{Hom}_B(X,N)\bigr).
    \]
\end{enumerate}
\end{theorem}




\begin{theorem}
    The following three conditions on $V_A$ are equivalent.
    \begin{enumerate}
        \item[(1)] $V$ is a Noetherian module.
        \item[(2)] (Ascending chain condition.) For every infinite chain of $A$-submodules of $V$,
        \begin{equation}\tag{2.1}
            V_1 \subset V_2 \subset \cdots \subset V_n \subset \cdots,
        \end{equation}
        there is a number $m$ such that $V_m = V_{m+1} = \cdots$.
        \item[(3)] Every $A$-submodule of $V$ is finitely generated over $A$.
    \end{enumerate}
\end{theorem}

\begin{remark}[Finiteness in Category Theory]
    \textbf{(1) Finite Presentation via Colimits:} \\
    In classical module categories, ``finiteness'' is characterized by filtered colimits. A module $M$ is finitely presented if and only if the functor $\operatorname{Hom}(M, -)$ commutes with filtered colimits:
    \[
        \varinjlim \operatorname{Hom}(M, N_i) \xrightarrow{\sim} \operatorname{Hom}(M, \varinjlim N_i)
    \]
    Intuitively, any morphism from $M$ to an infinite filtered system must factor through a finite stage.

    \vspace{1ex}
    \noindent \textbf{(2) Finiteness via Functor Preservation in $\infty$-Categories:} \\
    In higher category theory, objects lack ``elements.'' Instead, an object $X$ is defined as ``finite'' based on how its associated functors interact with limits and colimits.
    \begin{itemize}
        \item \textbf{Compact (Finitely Presented) Objects:} \\
        The functor $\operatorname{Map}(X, -)$ commutes with \textit{filtered colimits}. \\
        \textit{Example:} Finite CW complexes in the $\infty$-category of Spaces.

        \item \textbf{Dualizable Objects:} \\
        The functor $X \otimes -$ commutes with \textit{all limits} (e.g., infinite products). This happens because having a dual allows the left adjoint $X \otimes -$ to act like the right adjoint $\underline{\operatorname{Map}}(X^\vee, -)$. \\
        \textit{Example:} Finite spectra in the stable $\infty$-category of Spectra.

        \item \textbf{Perfect Complexes:} \\
        The ultimate intersection of both finiteness conditions. $\mathbb{R}\operatorname{Map}(X, -)$ preserves filtered colimits, and $X \otimes^\mathbb{L} -$ preserves limits. \\
        \textit{Example:} Perfect complexes in the derived $\infty$-category $\mathcal{D}(R)$ or $\operatorname{QCoh}(X)$ in derived algebraic geometry.
    \end{itemize}
\end{remark}

\begin{theorem}
    The following three conditions on a module $V_A$ are equivalent:
    \begin{enumerate}
        \item[(1)] $V$ is an Artinian module.
        \item[(2)] (Descending chain condition.) For every infinite chain of $A$-submodules of $V$,
        \[
            V_1 \supset V_2 \supset \cdots \supset V_n \supset \cdots,
        \]
        there is a number $m$ such that $V_m = V_{m+1} = \cdots$.
        \item[(3)] (Minimal condition.) Every non-empty set of $A$-submodules of $V$ contains a minimal element with respect to inclusion.
    \end{enumerate}
\end{theorem}

\begin{theorem}\label{thm:na-critria}
Let $W$ be an $A$-submodule of $V$. Then $V$ is Noetherian (respectively, Artinian) if and only if both $W$ and $V/W$ are Noetherian (respectively, Artinian).
\end{theorem}

\begin{definition}
A nonzero $A$-module $V$ is said to be \textbf{irreducible} (or \textbf{simple}) if it has no proper $A$-submodule other than $0$. 

A chain of $A$-submodules of $V$
$$V = V_0 \supset V_1 \supset \cdots \supset V_n = 0$$
is said to be a \textbf{composition series} of $V$ if each $V_i/V_{i+1}$ is irreducible. And the simple modules
$$V/V_1, V_1/V_2, \dots, V_{n-1}/V_n = V_{n-1}$$
are called the \textbf{composition factors} of the series.
\end{definition}

\begin{theorem}
A non-zero module has a composition series if and only if it is both Noetherian and Artinian.
\end{theorem}
\begin{proof}
    Use \autoref{thm:na-critria} and induction on the length of chain.
\end{proof}

\begin{theorem}[Jordan--Hölder]
If \(V\) has a composition series, any two composition factors obtained
from the two series of \(V\) are the same up to order of occurrence and
isomorphism.
\end{theorem}



\begin{theorem}[Equivalent descriptions of the Jacobson radical]\label{thm:eq des of J}
Let \(A\) be a unital associative ring. For \(a\in A\), the following
conditions are equivalent:
\begin{enumerate}
    \item[(i)] \(a\in J(A)\).
    \item[(ii)] \(a\) annihilates every simple left \(A\)-module, i.e.
    \(aS=0\) for every simple left \(A\)-module \(S\).
    \item[(iii)] \(a\) annihilates every simple right \(A\)-module, i.e.
    \(Ta=0\) for every simple right \(A\)-module \(T\).
    \item[(iv)] \(a\) belongs to every maximal left ideal of \(A\).
    \item[(v)] \(a\) belongs to every maximal right ideal of \(A\).
    \item[(vi)] \(1-xa\in A^\times\) for every \(x\in A\).
    \item[(vii)] \(1-ax\in A^\times\) for every \(x\in A\).
\end{enumerate}
Equivalently,
\[
J(A)=\bigcap_{S\ \mathrm{simple}}\operatorname{Ann}_A(S)
     =\bigcap_{\mathfrak m\ \mathrm{maximal\ left}}\mathfrak m
     =\bigcap_{\mathfrak n\ \mathrm{maximal\ right}}\mathfrak n .
\]
Moreover, \(J(A)\) is the largest two-sided ideal \(I\subset A\) such that
\(1-a\in A^\times\) for every \(a\in I\).
\end{theorem}

\begin{theorem}
Let \(A\) be a unital associative ring. If \(I\subset A\) is a nil left(right, double-sided) ideal,
then \(I\subset J(A)\).
\end{theorem}
\begin{proof}
    Use the (vi) of \autoref{thm:eq des of J}.
\end{proof}

\begin{theorem}[Jacobson radical of an Artinian ring] \label{thm:jacobson-radical-artinian} 
    Let \(A\) be a unital associative ring. If \(A\) is left Artinian or right Artinian, then \(J(A)\) is the largest nilpotent two-sided ideal of \(A\).
\end{theorem}

\begin{theorem} \label{thm:simple-left-modules-quotients} 
    Let \(A\) be a unital associative ring. If \(S\) is a simple left \(A\)-module, then there exists a maximal left ideal \(\mathfrak m\subset A\) such that \[ S\cong A/\mathfrak m \] as left \(A\)-modules. 
\end{theorem}

\begin{remark} \label{rem:annihilator-of-simple-quotient} 
    Let \(A\) be a unital associative ring, and let \(\mathfrak m\subset A\) be a maximal left ideal. Then \(A/\mathfrak m\) is a simple left \(A\)-module, and \[ \operatorname{Ann}_A(A/\mathfrak m) = \{a\in A\mid aA\subset \mathfrak m\}. \] This is the largest two-sided ideal contained in \(\mathfrak m\). In general it need not equal \(\mathfrak m\), since \(\mathfrak m\) is only a left ideal. Thus \(\mathfrak m\) is the annihilator of the element \(1+\mathfrak m\), whereas \(\operatorname{Ann}_A(A/\mathfrak m)\) is the annihilator of the whole module. 
\end{remark}

\begin{theorem}[Azumaya--Nakayama] \label{thm:azumaya-nakayama} 
    Let \(A\) be a unital associative ring. 
    \begin{enumerate} 
        \item[(i)] Let \(V\) be a left \(A\)-module and let \(W\) be an \(A\)-submodule of \(V\) such that \(V/W\) is finitely generated over \(A\). Then \[ V=W+J(A)V\Longrightarrow V=W. \] 
        \item[(ii)] If \(V\) is a finitely generated left \(A\)-module, then \[ J(A)V=V\Longrightarrow V=0. \] 
    \end{enumerate} 
\end{theorem}

\begin{theorem}
\label{thm:finite-length-semisimple-radical}
Let \(A\) be a unital associative ring, and let \(M\) be a left \(A\)-module
of finite length. Then
\[
M \text{ is semisimple}
\Longleftrightarrow
\operatorname{rad}(M)=0.
\]
\end{theorem}

\begin{theorem} \label{thm:semisimple-and-jacobson-radical} 
    Let \(A\) be a unital associative ring. 
    \begin{enumerate} 
        \item[(i)] If \(A\) is semisimple, then \(J(A)=0\). 
        \item[(ii)] If \(A\) is left Artinian or right Artinian, then \[ A \text{ is semisimple} \Longleftrightarrow J(A)=0. \] 
    \end{enumerate} 
\end{theorem}


\begin{definition}
\label{def:radical-of-module}
Let \(A\) be a unital associative ring, and let \(M\) be a left
\(A\)-module. The radical of \(M\) is defined by
\[
\operatorname{rad}(M)
=
\bigcap_{N\subset M\ \mathrm{maximal}} N,
\]
where the intersection is taken over all maximal submodules of \(M\).
\end{definition}

\begin{corollary}
\label{cor:jacobson-radical-contained-in-module-radical}
Let \(A\) be a unital associative ring, and let \(M\) be a left \(A\)-module.
Then
\[
J(A)M\subset \operatorname{rad}(M).
\]
\end{corollary}

\begin{theorem} \label{thm:semisimple-module-radical} 
    Let \(A\) be a unital associative ring, and let \(M\) be a left \(A\)-module. If \(M\) is semisimple, then \[ \operatorname{rad}(M)=0. \] 
\end{theorem}


\begin{theorem}
\label{thm:module-radical-artinian}
Let \(A\) be a left Artinian ring, and let \(M\) be a finitely generated left
\(A\)-module. Then
\[
\operatorname{rad}(M)=J(A)M.
\]
\end{theorem}

\begin{proof}
    Use \autoref{thm:semisimple-module-radical}.
\end{proof}



\begin{theorem}
\label{thm:direct-sum-and-idempotents}
Let \(A\) be a unital associative ring, and let \(M\) be a left
\(A\)-module. A decomposition
\[
M=M_1\oplus\cdots\oplus M_n
\]
is equivalent to a decomposition of \(\operatorname{id}_M\) in
\(\operatorname{End}_A(M)\) as
\[
\operatorname{id}_M=p_1+\cdots+p_n,
\]
where \(p_i^2=p_i\) and \(p_ip_j=0\) for \(i\neq j\). Under this
correspondence,
\[
M_i=\operatorname{Im}(p_i).
\]
In particular, \(M\) is indecomposable if and only if \(\operatorname{id}_M\)
is a primitive idempotent in \(\operatorname{End}_A(M)\).
\end{theorem}


\begin{theorem}
\label{thm:endomorphism-ring-of-regular-module}
Let \(A\) be a unital associative ring. Then
\[
\operatorname{End}_A(A_A)\cong A \quad \operatorname{End}_A(_A A)\cong A^{\text{op}}.
\]
\end{theorem}


\begin{theorem} \label{thm:idempotents-and-ideal-decompositions} 
    Let \(A\) be a unital associative ring, and let \(e\in A\) be an idempotent. 
    \begin{enumerate} 
        \item[(i)] If \(e=e_1+\cdots+e_n\) is a decomposition into pairwise orthogonal idempotents, then \[ eA=e_1A\oplus\cdots\oplus e_nA \] as right \(A\)-modules, and \[ Ae=Ae_1\oplus\cdots\oplus Ae_n \] as left \(A\)-modules. 
        \item[(ii)] Conversely, every direct sum decomposition \[ eA=I_1\oplus\cdots\oplus I_n \] into right ideals is of the form \(I_i=e_iA\), where \(e=e_1+\cdots+e_n\) is a decomposition into pairwise orthogonal idempotents. Similarly, every direct sum decomposition \[ Ae=L_1\oplus\cdots\oplus L_n \] into left ideals is of the form \(L_i=Ae_i\), where \(e=e_1+\cdots+e_n\) is a decomposition into pairwise orthogonal idempotents. 
    \end{enumerate} 
\end{theorem}


\begin{theorem}
\label{thm:primitive-idempotents-for-modules}
Let \(A\) be a unital associative ring, let \(M\) be a left \(A\)-module, and
set \(E=\operatorname{End}_A(M)\). For a nonzero idempotent \(p\in E\), write
\(pM=\operatorname{Im}(p)\). Then
\[
p \text{ is primitive in } E
\Longleftrightarrow
pM \text{ is indecomposable as a left } A\text{-module}.
\]
Moreover,
\[
\operatorname{End}_A(pM)\cong pEp .
\]
\end{theorem}


\begin{theorem}
\label{thm:indecomposable-local-endomorphism-ring}
Let \(A\) be a unital associative ring, and let \(V\) be a left
\(A\)-module. Suppose that \(V\) has a composition series. Then
\[
V \text{ is indecomposable}
\Longleftrightarrow
\operatorname{End}_A(V) \text{ is a local ring}.
\]
\end{theorem}

\begin{theorem}[Krull--Schmidt--Azumaya]
\label{thm:krull-schmidt-azumaya}
Let \(A\) be a unital associative ring, and let \(V\) be a left
\(A\)-module. Suppose that
\[
V=V_1\oplus\cdots\oplus V_n
\]
is a finite direct sum decomposition into indecomposable submodules, and that
\(\operatorname{End}_A(V_i)\) is a local ring for every \(i\).

Then this decomposition is unique up to isomorphism and permutation. That is,
if
\[
V=W_1\oplus\cdots\oplus W_m
\]
is another finite direct sum decomposition into indecomposable submodules,
then \(m=n\), and there exists a permutation \(\sigma\in S_n\) such that
\[
V_i\cong W_{\sigma(i)}
\]
for all \(1\leq i\leq n\).
\end{theorem}



\begin{lemma}
\label{lem:idempotent-finite-decomposition-artinian-module}
Let \(A\) be a unital associative ring, and let \(M\) be an Artinian left
\(A\)-module. Set \(E=\operatorname{End}_A(M)\). Then every nonzero
idempotent \(p\in E\) admits a finite decomposition
\[
p=p_1+\cdots+p_n
\]
where \(p_1,\dots,p_n\in E\) are pairwise orthogonal primitive idempotents.
Equivalently,
\[
\operatorname{Im}(p)
=
\operatorname{Im}(p_1)\oplus\cdots\oplus \operatorname{Im}(p_n)
\]
is a finite direct sum decomposition into indecomposable left
\(A\)-modules.
\end{lemma}

\begin{proof}
Since \(p\) is idempotent, \(\operatorname{Im}(p)\) is a direct summand of
\(M\). Hence \(\operatorname{Im}(p)\) is Artinian. Repeatedly decompose
\(\operatorname{Im}(p)\) whenever it is decomposable. This process terminates,
because otherwise one obtains an infinite strictly descending chain of
submodules of \(M\). Thus
\[
\operatorname{Im}(p)=M_1\oplus\cdots\oplus M_n
\]
with each \(M_i\) indecomposable.

Let \(p_i\) be the projection from \(M\) onto \(M_i\), with kernel equal to
the sum of the other summands and \(\ker(p)\). Then the \(p_i\) are pairwise
orthogonal idempotents in \(E\), and
\[
p=p_1+\cdots+p_n.
\]
Since \(\operatorname{Im}(p_i)=M_i\) is indecomposable, \(p_i\) is primitive.
\end{proof}





\begin{corollary}
\label{cor:krull-schmidt-artinian}
Let \(A\) be a left Artinian ring, and let \(M\) be a finitely generated
left \(A\)-module. Then \(M\) admits a finite direct sum decomposition
\[
M=M_1\oplus\cdots\oplus M_n
\]
where each \(M_i\) is indecomposable. Moreover, this decomposition is unique
up to isomorphism and permutation.
\end{corollary}

\begin{proof}
    Use \autoref{lem:idempotent-finite-decomposition-artinian-module} and \autoref{thm:krull-schmidt-azumaya}.
\end{proof}


\begin{theorem}
\label{thm:projective-module-characterizations}
Let \(A\) be a unital associative ring, and let \(P\) be a left
\(A\)-module. The following conditions are equivalent:
\begin{enumerate}
    \item[(i)] \(P\) is projective.

    \item[(ii)] For every epimorphism \(M\twoheadrightarrow N\) and every
    homomorphism \(P\to N\), there exists a lift \(P\to M\).

    \item[(iii)] The functor
    \[
    \operatorname{Hom}_A(P,-):A\text{-}\mathrm{Mod}\longrightarrow
    \mathrm{Ab}
    \]
    is exact.

    \item[(iv)] \(P\) is a direct summand of a free left \(A\)-module.

    \item[(v)] Every epimorphism \(M\twoheadrightarrow P\) splits.
\end{enumerate}
If \(P\) is finitely generated, then these conditions are also equivalent to
\(P\) being a direct summand of \(A^n\) for some \(n\geq 1\).
\end{theorem}

\begin{theorem}
\label{thm:injective-module-characterizations}
Let \(A\) be a unital ring, and let \(I\) be a left \(A\)-module. The following
conditions are equivalent:
\begin{enumerate}
    \item[(i)] \(I\) is injective, i.e., for every submodule \(M\subset N\) and
    every homomorphism \(f:M\to I\), there exists an extension \(g:N\to I\)
    such that \(g|_M=f\).

    \item[(ii)] The functor
    \[
    \operatorname{Hom}_A(-,I):(A\text{-Mod})^{\mathrm{op}}\longrightarrow
    \mathrm{Ab}
    \]
    is exact.

    \item[(iii)] (Baer criterion) For every left ideal \(J\subset A\) and
    homomorphism \(f:J\to I\), there exists an extension \(g:A\to I\)
    with \(g|_J=f\).
\end{enumerate}
\end{theorem}


\begin{theorem}
\label{thm:flat-module-characterizations}
Let \(A\) be a unital associative ring, and let \(F\) be a left
\(A\)-module. The following conditions are equivalent:
\begin{enumerate}
    \item[(i)] \(F\) is flat, i.e. the functor
    \[
    -\otimes_A F:\mathrm{Mod}\text{-}A\longrightarrow \mathrm{Ab}
    \]
    is exact. Equivalently, for every monomorphism \(M'\hookrightarrow M\)
    of right \(A\)-modules, the induced map
    \[
    M'\otimes_A F\longrightarrow M\otimes_A F
    \]
    is a monomorphism.

    \item[(ii)] For every finitely generated right ideal \(I\subset A\), the
    natural map
    \[
    I\otimes_A F\longrightarrow A\otimes_A F
    \]
    is injective.

    \item[(iii)] \(F\) is a filtered colimit of finitely generated free left
    \(A\)-modules.
\end{enumerate}
\end{theorem}


\begin{theorem}[Wedderburn--Artin]
\label{thm:wedderburn-left-module}
Let \(A\) be a simple Artinian ring. Then there exist a division ring \(D\) and
a positive integer \(n\) such that
\[
A \cong M_n(D),
\]
the ring of \(n\times n\) matrices over \(D\).

Moreover, as a left \(A\)-module, the left regular module \({}_A A\) decomposes as
\[
{}_A A \cong S^{\oplus n},
\]
where \(S = Ae_{11}\) is a minimal left ideal of \(A\), simple and projective,
corresponding to the first column of the matrix block. Each of the \(n\) summands
is isomorphic to \(S\), so the left regular module is a direct sum of
mutually isomorphic minimal left ideals.
\end{theorem}



\begin{theorem}[Idempotent lifting modulo a nilpotent ideal]
\label{thm:idempotent-lifting-nilpotent-ideal}
Let \(A\) be a unital associative ring, and let \(I\subset A\) be a
two-sided nilpotent ideal. Let \(\pi:A\to A/I\) be the quotient map.

\begin{enumerate}
    \item[(i)] Every idempotent of \(A/I\) lifts to an idempotent of \(A\).
    That is, if \(\bar e\in A/I\) satisfies \(\bar e^2=\bar e\), then there
    exists an idempotent \(e\in A\) such that \(\pi(e)=\bar e\).

    \item[(ii)] More generally, if
    \[
    \bar e=\bar e_1+\cdots+\bar e_n
    \]
    is a decomposition into pairwise orthogonal idempotents in \(A/I\), and
    \(e\in A\) is an idempotent lift of \(\bar e\), then there exist pairwise
    orthogonal idempotents \(e_1,\dots,e_n\in A\) such that
    \[
    e=e_1+\cdots+e_n,
    \qquad
    \pi(e_i)=\bar e_i
    \]
    for all \(1\leq i\leq n\).
\end{enumerate}
\end{theorem}


\begin{theorem} \label{thm:primitive-idempotent-lifting} 
    Let \(A\) be a unital associative ring, and let \(I\subset A\) be a two-sided nilpotent ideal. 
    Set \(\overline{A}=A/I\), and let \(\pi:A\to \overline{A}\) be the quotient map. If \(e\in A\) is an idempotent and \(\overline e=\pi(e)\), then \[ e \text{ is primitive in } A \Longleftrightarrow \overline e \text{ is primitive in } \overline A . \] 
\end{theorem}

\begin{lemma} \label{lem:essential-epimorphism-nakayama} 
    Let \(A\) be a unital associative ring, and let \(f:M\to N\) be an epimorphism of left \(A\)-modules. If \(M\) is finitely generated and \[ \ker f\subset J(A)M, \] then \(f\) is essential. That is, if \(L\subset M\) is a submodule such that \(f(L)=N\), then \(L=M\). 
\end{lemma}


\begin{theorem}
\label{thm:projective-covers-over-semiperfect-rings}
Let \(A\) be a semiperfect ring, let \(J=J(A)\), and set
\(\overline A=A/J\).

\begin{enumerate}
    \item[(i)] Every finitely generated left \(A\)-module \(M\) has a
    projective cover
    \[
    p:P\twoheadrightarrow M,
    \]
    where \(P\) is a finitely generated projective left \(A\)-module. Moreover,
    the induced map
    \[
    P/JP\longrightarrow M/JM
    \]
    is an isomorphism.

    \item[(ii)] Let \(f:P\twoheadrightarrow M\) be an epimorphism of left
    \(A\)-modules, and suppose that \(P\) is finitely generated and
    projective. Then
    \[
    f \text{ is essential}
    \Longleftrightarrow
    \ker f\subset JP.
    \]

    \item[(iii)] If \(p:P\twoheadrightarrow M\) is the projective cover in
    \textup{(i)}, then the induced map
    \[
    \overline p:P/JP\longrightarrow M/JM
    \]
    is a projective cover in \(\overline A\text{-}\mathrm{Mod}\).
\end{enumerate}
\end{theorem}

\section{Representation Theory}

\begin{theorem}[Characters are class functions]
\label{thm:characters-are-class-functions}
Let \(G\) be a finite group, and let \(\chi\) be a complex character of
\(G\). If \(g,h\in G\) are conjugate, then
\[
\chi(g)=\chi(h).
\]
Equivalently, every character of \(G\) is constant on conjugacy classes.
\end{theorem}

\begin{theorem}[Number of irreducible characters]
\label{thm:number-of-irreducible-characters}
Let \(G\) be a finite group, and let \(\operatorname{Irr}(G)\) be the set of
irreducible complex characters of \(G\). Then
\[
|\operatorname{Irr}(G)|
=
|\operatorname{Cl}(G)|,
\]
where \(\operatorname{Cl}(G)\) denotes the set of conjugacy classes of \(G\).
\end{theorem}

\begin{theorem}[Linear characters and the abelianization]
\label{thm:linear-characters-and-abelianization}
Let \(G\) be a finite group, and let \(G'=[G,G]\) be its commutator subgroup.
Then the one-dimensional complex characters of \(G\) are exactly the
characters inflated from the abelianization \(G/G'\). Equivalently,
\[
\operatorname{Irr}_1(G)
\cong
\operatorname{Hom}(G/G',\mathbb C^\times).
\]
In particular,
\[
|\operatorname{Irr}_1(G)|=|G/G'|.
\]
\end{theorem}

\begin{theorem}[Degree formula for irreducible characters]
\label{thm:degree-formula-for-irreducible-characters}
Let \(G\) be a finite group, and let \(\operatorname{Irr}(G)\) be the set of
irreducible complex characters of \(G\). Then
\[
|G|
=
\sum_{\chi\in\operatorname{Irr}(G)}\chi(1)^2.
\]
Moreover, for every \(\chi\in\operatorname{Irr}(G)\),
\[
\chi(1)\mid |G|.
\]
\end{theorem}


\begin{theorem}[Row orthogonality of irreducible characters]
\label{thm:row-orthogonality-irreducible-characters}
Let \(G\) be a finite group, and let \(\operatorname{Irr}(G)\) be the set of
irreducible complex characters of \(G\). For class functions \(\chi,\psi\),
define
\[
\langle \chi,\psi\rangle
=
\frac{1}{|G|}
\sum_{g\in G}\chi(g)\overline{\psi(g)}.
\]
Then for \(\chi,\psi\in\operatorname{Irr}(G)\),
\[
\langle \chi,\psi\rangle=\delta_{\chi,\psi}.
\]
Equivalently, if \(C_1,\dots,C_r\) are the conjugacy classes of \(G\), then
\[
\frac{1}{|G|}
\sum_{\ell=1}^r
|C_\ell|\chi(C_\ell)\overline{\psi(C_\ell)}
=
\delta_{\chi,\psi}.
\]
\end{theorem}

\begin{theorem}[Column orthogonality of irreducible characters]
\label{thm:column-orthogonality-irreducible-characters}
Let \(G\) be a finite group, and let \(\operatorname{Irr}(G)\) be the set of
irreducible complex characters of \(G\). If \(g,h\in G\), then
\[
\sum_{\chi\in\operatorname{Irr}(G)}
\chi(g)\overline{\chi(h)}
=
\begin{cases}
|C_G(g)|, & g \text{ is conjugate to } h,\\
0, & g \text{ is not conjugate to } h.
\end{cases}
\]
In particular,
\[
\sum_{\chi\in\operatorname{Irr}(G)}|\chi(g)|^2
=
|C_G(g)|.
\]
\end{theorem}



\begin{theorem}[Basic facts on ordinary characters of finite groups]
\label{thm:basic-character-theory}
Let \(G\) be a finite group, and let \(\operatorname{Irr}(G)\) denote the set
of irreducible complex characters of \(G\). Then the following facts hold.

\begin{enumerate}
    \item[(i)] Every character is a class function. That is, if \(g,h\in G\)
    are conjugate, then
    \[
    \chi(g)=\chi(h).
    \]

    \item[(ii)] The irreducible characters form an orthonormal basis of the
    space of complex class functions on \(G\), with respect to the inner
    product
    \[
    \langle \chi,\psi\rangle
    =
    \frac{1}{|G|}\sum_{g\in G}\chi(g)\overline{\psi(g)}.
    \]
    In particular,
    \[
    \langle \chi,\psi\rangle=\delta_{\chi,\psi}
    \]
    for \(\chi,\psi\in \operatorname{Irr}(G)\).

    \item[(iii)] The number of irreducible characters equals the number of
    conjugacy classes of \(G\):
    \[
    |\operatorname{Irr}(G)|
    =
    |\operatorname{Cl}(G)|.
    \]

    \item[(iv)] The one-dimensional characters of \(G\) are exactly the
    characters inflated from the abelianization \(G/G'\), where
    \(G'=[G,G]\). Equivalently,
    \[
    \operatorname{Irr}_1(G)
    \cong
    \operatorname{Hom}(G/G',\mathbb C^\times).
    \]
    In particular,
    \[
    |\operatorname{Irr}_1(G)|=|G/G'|.
    \]

    \item[(v)] The regular character \(\rho_{\mathrm{reg}}\) satisfies
    \[
    \rho_{\mathrm{reg}}(1)=|G|,
    \qquad
    \rho_{\mathrm{reg}}(g)=0
    \quad (g\neq 1),
    \]
    and decomposes as
    \[
    \rho_{\mathrm{reg}}
    =
    \sum_{\chi\in\operatorname{Irr}(G)}\chi(1)\chi.
    \]

    \item[(vi)] The degrees of the irreducible characters satisfy
    \[
    |G|
    =
    \sum_{\chi\in\operatorname{Irr}(G)}\chi(1)^2.
    \]
    Moreover, for every \(\chi\in\operatorname{Irr}(G)\),
    \[
    \chi(1)\mid |G|.
    \]

    \item[(vii)] The row orthogonality relations are
    \[
    \frac{1}{|G|}
    \sum_{C\in\operatorname{Cl}(G)}
    |C|\chi(C)\overline{\psi(C)}
    =
    \delta_{\chi,\psi}.
    \]

    \item[(viii)] The column orthogonality relations are as follows. If
    \(g,h\in G\), then
    \[
    \sum_{\chi\in\operatorname{Irr}(G)}
    \chi(g)\overline{\chi(h)}
    =
    \begin{cases}
    |C_G(g)|, & g \text{ is conjugate to } h,\\
    0, & g \text{ is not conjugate to } h.
    \end{cases}
    \]
    In particular,
    \[
    \sum_{\chi\in\operatorname{Irr}(G)}|\chi(g)|^2
    =
    |C_G(g)|.
    \]

    \item[(ix)] If \(z\in Z(G)\) and \(\chi\in\operatorname{Irr}(G)\), then
    \(z\) acts by a scalar on any representation affording \(\chi\). Hence
    \[
    |\chi(z)|=\chi(1).
    \]
    Consequently,
    \[
    \chi(1)^2\leq |G:Z(G)|.
    \]

    \item[(x)] If \(N\triangleleft G\), then inflation gives an injection
    \[
    \operatorname{Irr}(G/N)\hookrightarrow \operatorname{Irr}(G).
    \]
    Its image consists exactly of those irreducible characters \(\chi\) of
    \(G\) such that
    \[
    N\subset \ker \chi.
    \]

    \item[(xi)] For every character \(\chi\), its kernel is
    \[
    \ker\chi
    =
    \{g\in G\mid \chi(g)=\chi(1)\}.
    \]
    If \(\chi\) is irreducible, this is the kernel of any representation
    affording \(\chi\), and hence is a normal subgroup of \(G\).

    \item[(xii)] If \(H\leq G\), then induction and restriction satisfy
    Frobenius reciprocity:
    \[
    \langle \operatorname{Ind}_H^G\varphi,\chi\rangle_G
    =
    \langle \varphi,\operatorname{Res}_H^G\chi\rangle_H
    \]
    for every character \(\varphi\) of \(H\) and every character \(\chi\) of
    \(G\).

    \item[(xiii)] If \(G\) acts on a finite set \(X\), then the associated
    permutation character is
    \[
    \pi_X(g)=|\operatorname{Fix}_X(g)|.
    \]

    \item[(xiv)] If \(\chi\) and \(\psi\) are characters of \(G\), then
    \(\chi+\psi\) is the character of a direct sum and \(\chi\psi\) is the
    character of a tensor product. The character of the dual representation is
    \[
    \chi^\vee(g)=\chi(g^{-1})=\overline{\chi(g)}.
    \]

    \item[(xv)] For every character \(\chi\) and every \(g\in G\), the value
    \(\chi(g)\) is an algebraic integer. In particular, if \(\chi(g)\in
    \mathbb Q\), then \(\chi(g)\in\mathbb Z\).

    \item[(xvi)] If \(A\triangleleft G\) is abelian, then for every
    \(\chi\in\operatorname{Irr}(G)\),
    \[
    \chi(1)\mid |G:A|.
    \]
\end{enumerate}
\end{theorem}




























